% TikZ 图：有界磁场圆弧
% Day 22：洛伦兹力
\begin{tikzpicture}[scale=1.22,>=Latex,line cap=round,line join=round]
  \path[use as bounding box] (-1.75,-2.05) rectangle (3.65,2.00);
  \draw[thick,fill=blue!4] (0,-1.48) rectangle (2.75,1.48);
  \foreach \x in {0.45,1.10,1.75,2.40} {
    \foreach \y in {-1.02,-0.34,0.34,1.02} {
      \draw[blue!70!black,thick] (\x-0.055,\y-0.055) -- (\x+0.055,\y+0.055);
      \draw[blue!70!black,thick] (\x-0.055,\y+0.055) -- (\x+0.055,\y-0.055);
    }
  }
  \node[blue!70!black,font=\small] at (2.13,1.13) {$B$ 向里};

  \coordinate (O) at (0,0.35);
  \coordinate (A) at (0,-1.05);
  \coordinate (C) at (1.21,-0.35);
  \draw[->,red,very thick] (-1.30,-1.05) -- (A);
  \node[red,below] at (-0.68,-1.05) {$v_1$};
  \draw[red,very thick,->] (A) arc (-90:-30:1.40);
  \draw[->,red,very thick] (C) -- ++(0.72,1.25);
  \node[red,right] at (1.70,0.58) {$v_2$};

  \fill[red] (A) circle (2pt);
  \node[red,below left=2pt] at (A) {$A$};
  \node[red] at (0.35,-1.30) {$+q$};
  \fill[red] (C) circle (2pt);
  \node[red,below right=2pt] at (C) {$C$};
  \fill (O) circle (1.6pt) node[left=3pt] {$O$};
  \draw[gray,dashed,thick] (O) -- (A);
  \node[gray,left] at (0,-0.35) {$r$};
  \draw[gray,dashed,thick] (O) -- (C);
  \node[gray,above right] at (0.58,0.02) {$r$};
  \node[font=\small] at (1.35,-1.78)
    {入射、出射速度的垂线相交于圆心};
\end{tikzpicture}
